% TOP_PROBLEM: {"id": 8100008, "title": "Rigorous Diffraction from Two Slits", "category_id": 16, "category": {"id": 16, "name": "physics", "display_name": "Mathematical Physics", "description": "Problems at the intersection of mathematics and physics.", "slug": "physics", "order_index": 16, "created_at": "2026-07-31T15:26:25.671Z"}, "status": "open", "statusQualification": "The source does not fix slit geometry, screen boundary condition or incident field in its short formulation; these must be stated for any proposed solution. The Helmholtz description is explanatory context, not an invented choice of Dirichlet or Neumann target."}
% FIRST_POSTED: 2026-10-11
% ENTRY_KIND: statement_only
% SUBMISSION: {"kind": "statement_only", "category": "top_problems", "problem_id": "8100008", "proposal_type": "established", "source": "Percy Deift, Some Open Problems in Random Matrix Theory and the Theory of Integrable Systems. II, SIGMA 13 (2017), 016, final two-slit question. https://arxiv.org/abs/1703.04931; Primary full-text formulation and 2017 historical progress discussion. https://arxiv.org/html/1703.04931v1", "status_evidence": "Inherited source-catalogue status open preserved. Primary source accessed 11 October 2026; compared against the supplied catalogue excerpt. Historical wording and partial-result remarks were checked for scope; no comprehensive current-literature or proof-verification audit was performed.", "status_checked_on": "2026-10-11", "duplicate_check": "Full 15,773-record source registry checked for identical normalized statements; available site registry checked for identical normalized title/statement and canonical-ID collision. Canonical ID is absent from the site registry. No mathematical-equivalence literature audit performed.", "attribution": "Percy Deift. Statement preparation: OpenAI GPT-6 Codex assistant, generation 11 October 2026; source-matching and notation editing by the assistant, no independent human mathematical verification. The consulted arXiv HTML version states CC BY-SA 4.0; preserve applicable attribution and share-alike provenance. Inherited rights/provenance note: Public arXiv source TeX; arXiv reuse license NEEDS_REVIEW"}
% SOURCE_STATEMENT: "Give a rigorous explicit solution of the fixed-frequency scalar-wave diffraction problem for two finite slits in the plane, including asymptotics of the solution."
% Preparation and final self-audit: CONTRIBUTING.md prompts C and D; source/audit artifacts retained in the local physics_statements_2026_10_11 batch.
% Imported from: physics_statements.tex; UnsolvedMath ID 8100008
% !TeX program = lualatex
\documentclass[11pt]{article}
\usepackage[margin=25mm]{geometry}
\usepackage{fontspec}
\setmainfont{Latin Modern Roman}
\usepackage{amsmath,amssymb,mathtools}
\usepackage{unicode-math}
\setmathfont{Latin Modern Math}
\usepackage{xurl,hyperref}
\hypersetup{colorlinks=true,urlcolor=blue}
\setcounter{secnumdepth}{0}
\setlength{\parindent}{0pt}
\setlength{\parskip}{5pt}
\setlength{\emergencystretch}{3em}
\newcommand{\sref}[2]{\hyperlink{source:#1:#2}{[S#2]}}
\newcommand{\cataloguescope}[1]{\paragraph{Recorded target.}#1}
\newcommand{\R}{\mathbb R}
\newcommand{\C}{\mathbb C}
\newcommand{\N}{\mathbb N}
\newcommand{\Z}{\mathbb Z}
\newcommand{\Q}{\mathbb Q}
\newcommand{\F}{\mathbb F}
\newcommand{\eps}{\varepsilon}
\begin{document}
% UnsolvedMath ID 8100008; source catalogue code AMR-080-0008
\setcounter{section}{8100007}
\section{Rigorous Diffraction from Two Slits}\label{8100008}
{\small UnsolvedMath ID 8100008 · Mathematical Physics}
\subsection{Definitions and mathematical statement}

\paragraph{Shared notation.}$\mathbb N=\{1,2,\ldots\}$; the standard number fields have their usual meanings. All additional conventions are stated locally.


\paragraph{Definitions and notation.}At a fixed scalar-wave frequency the spatial field satisfies a Helmholtz equation $(\Delta+k^2)u=0$ away from the diffracting screen, with $k>0$. A diffraction problem includes an incident wave, boundary conditions on the screen and an outgoing radiation condition. Two finite slits introduce multiple edges and interactions absent from the single-slit setting.

\paragraph{Catalogue wording retained for traceability.}
Give a rigorous explicit solution of the fixed-frequency scalar-wave diffraction problem for two finite slits in the plane, including asymptotics of the solution.
\paragraph{Scope and source qualification.}The source does not fix slit geometry, screen boundary condition or incident field in its short formulation; these must be stated for any proposed solution. The Helmholtz description is explanatory context, not an invented choice of Dirichlet or Neumann target.
\subsection{Short English statement}
Find an explicit rigorous solution and asymptotic analysis for diffraction through two finite planar slits.
\subsection{Sources}
\begin{itemize}
\item[\hypertarget{source:8100008:1}{S1}] ULAM.AI and the UnsolvedMath Contributors, UnsolvedMath, record 8100008 (AMR-080-0008). Catalogue curation under CC BY 4.0; definitions and English exposition added locally.
\url{https://huggingface.co/datasets/ulamai/UnsolvedMath}.
\item[\hypertarget{source:8100008:2}{S2}] Percy Deift, Some Open Problems in Random Matrix Theory and the Theory of Integrable Systems. II, SIGMA 13 (2017), 016, final two-slit question.
\url{https://arxiv.org/abs/1703.04931}.
\item[\hypertarget{source:8100008:3}{S3}] Primary full-text formulation and 2017 historical progress discussion.
\url{https://arxiv.org/html/1703.04931v1}.
\end{itemize}
\label{end:8100008}
\end{document}
